\phantomsection\addcontentsline{toc}{section}{Space Engineering}
\ECchapter{09}{Space Engineering}{How do engineers design systems for extreme environments?}{ECDarkGrey}

\ECObjectives{
\begin{itemize}
\item Identify the main subsystems of an artificial satellite.
\item Explain how engineers manage launch loads, vacuum and temperature variations.
\item Propose a coherent architecture for a small Earth-observation mission.
\end{itemize}}

\begin{multicols}{2}
\begin{engineeringnews}
Reusable launch vehicles, commercial constellations and small satellites are changing access to
space. Lower launch costs make new missions possible, but every spacecraft must still survive an
extreme environment and operate for years without direct maintenance.
\end{engineeringnews}

\begin{ecarticle}
A spacecraft is designed around its \textbf{mission}. An Earth-observation satellite must carry a
camera or radar payload, point it accurately towards the ground, store the collected data and send
it to a ground station. Around this payload, engineers build a service platform containing the
structure, power supply, thermal control, communication system, onboard computer and attitude-control
system.

Launch is the first major challenge. The satellite experiences acceleration, vibration and intense
acoustic pressure. Its structure must therefore be light but stiff, while sensitive components are
mounted and tested so that they do not resonate dangerously. After separation from the launcher,
the environment changes completely: there is almost no atmosphere, heat cannot be removed by
convection and electronic components are exposed to radiation.

Electrical power usually comes from photovoltaic panels. Batteries supply the spacecraft during
eclipse, when Earth blocks the Sun. A power-control unit regulates voltage and distributes energy
to the payload, computer, heaters and transmitters. Because available power is limited, mission
operations are planned carefully: a high-power transmitter or instrument may only be activated at
specific times.

Thermal control is equally important. In sunlight, exposed surfaces may become very hot; in shadow,
they may cool rapidly. Engineers use insulation blankets, reflective coatings, heat pipes,
radiators and electrical heaters. The objective is not to keep every part at room temperature, but
to keep each component within its qualified operating range.

A satellite must also control its orientation, called \textbf{attitude}. Sun sensors, star trackers
and gyroscopes estimate its orientation. Reaction wheels rotate internally and make the spacecraft
turn in the opposite direction. Small thrusters can desaturate the wheels or perform orbital
corrections. Since repair is generally impossible, critical functions may be duplicated and the
onboard software must detect faults and switch to a safe mode.
\end{ecarticle}

\begin{tcolorbox}[title=\sffamily\bfseries Did You Know?,colback=yellow!10,colframe=ECOrange]
In low Earth orbit, a satellite travels at about \(7.8\,\mathrm{km\,s^{-1}}\) and completes one orbit
in roughly 90 minutes. It may therefore pass from sunlight to eclipse many times each day.
\end{tcolorbox}


\begin{engineeringfocus}
\textbf{Typical small-satellite architecture}

\begin{center}
\begin{tikzpicture}[font=\scriptsize,>=Latex,node distance=5mm and 5mm]
\tikzset{block/.style={draw=ECBlue,fill=ECLightBlue,rounded corners,align=center,
minimum width=21mm,minimum height=10mm}}
\node[block] (payload) {payload\\camera};
\node[block,right=of payload] (computer) {onboard\\computer};
\node[block,right=of computer] (radio) {radio and\\antenna};
\node[block,below=of payload] (adcs) {attitude\\control};
\node[block,below=of computer] (power) {solar array,\\battery, PCDU};
\node[block,below=of radio] (thermal) {thermal\\control};
\draw[->,thick,ECBlue] (payload) -- (computer);
\draw[<->,thick,ECBlue] (computer) -- (radio);
\draw[<->,thick,ECBlue] (computer) -- (adcs);
\draw[->,thick,ECOrange] (power) -- (payload);
\draw[->,thick,ECOrange] (power) -- (computer);
\draw[->,thick,ECOrange] (power) -- (radio);
\draw[<->,thick,ECGreen] (thermal) -- (computer);
\end{tikzpicture}
\end{center}

The payload performs the mission, while the platform keeps it powered, pointed, protected and
connected. Interfaces between subsystems are as important as the components themselves.
\end{engineeringfocus}

\begin{keyfigures}
\begin{tabularx}{\linewidth}{@{}Xr@{}}
Typical low Earth orbit altitude & 400--800 km\\
Orbital speed & about 7.8 km/s\\
One orbit & about 90 min\\
Vacuum pressure & extremely low\\
Repair after launch & usually impossible
\end{tabularx}
\end{keyfigures}

\begin{tcolorbox}[title=\sffamily\bfseries Thermal cycle in orbit,colback=white,colframe=ECDarkGrey]
\begin{center}
\begin{tikzpicture}
\begin{axis}[
width=\linewidth,height=46mm,
ymin=-60,ymax=100,
ylabel={Illustrative surface temperature (°C)},
symbolic x coords={Sunlight,Eclipse,Sunlight2,Eclipse2},
xtick=data,
xticklabels={Sunlight,Eclipse,Sunlight,Eclipse},
xticklabel style={font=\scriptsize},
yticklabel style={font=\scriptsize},
grid=major]
\addplot[very thick,mark=*] table[x=phase,y=temperature,col sep=comma]{data/EC09/orbit_temperature.csv};
\end{axis}
\end{tikzpicture}
\end{center}
\footnotesize Illustrative values only. Real temperatures depend on orbit, surface coating,
orientation, internal dissipation and thermal-control design.
\end{tcolorbox}

\begin{tcolorbox}[title=\sffamily\bfseries Three-axis attitude control,colback=white,colframe=ECBlue]
\begin{center}
\begin{tikzpicture}[scale=.9,font=\scriptsize,>=Latex]
\draw[fill=ECGrey,draw=ECDarkGrey,rounded corners] (-1.2,-.65) rectangle (1.2,.65);
\draw[fill=ECBlue!20,draw=ECBlue] (-2.8,-.45) rectangle (-1.25,.45);
\draw[fill=ECBlue!20,draw=ECBlue] (1.25,-.45) rectangle (2.8,.45);
\draw[->,very thick,ECRed] (0,0) -- (0,1.55) node[above]{yaw};
\draw[->,very thick,ECGreen] (0,0) -- (1.65,-1.05) node[right]{roll};
\draw[->,very thick,ECOrange] (0,0) -- (-1.65,-1.05) node[left]{pitch};
\node at (0,0) {satellite};
\end{tikzpicture}
\end{center}
Sensors estimate orientation; reaction wheels or thrusters create the required rotation around the
three axes.
\end{tcolorbox}

\begin{vocabulary}
\begin{tabularx}{\linewidth}{@{}lX@{}}
spacecraft & engin spatial\\
payload & charge utile\\
launch load & effort au lancement\\
solar array & panneau solaire\\
eclipse & passage dans l'ombre\\
attitude & orientation\\
reaction wheel & roue de réaction\\
star tracker & viseur d'étoiles\\
thruster & propulseur\\
redundancy & redondance
\end{tabularx}
\end{vocabulary}

\begin{questions}
\item Which mechanical loads affect a satellite during launch?
\item Why is convection unavailable in space?
\item What is the role of the battery during eclipse?
\item Name three devices used for attitude determination or control.
\item Why are critical functions sometimes duplicated?
\end{questions}

\begin{engineeringchallenge}
Design a 20 kg Earth-observation satellite. Draw a subsystem diagram and explain how the satellite
will generate and store power, control its temperature, point its camera, communicate with Earth
and react to one critical failure. Identify two design compromises involving mass, power, cost or
reliability.
\end{engineeringchallenge}

\begin{applicationexercise}
A satellite in low Earth orbit receives 950~W from its solar arrays in sunlight and
requires an average electrical power of 620~W. During a 36-minute eclipse, the battery
supplies the complete load. Estimate the minimum energy that must be delivered by the
battery in kWh. Then calculate the required nominal capacity if only 70\% of the battery
capacity may be used.
\end{applicationexercise}

\begin{speaking}
Present your satellite in a four-minute mission review. Explain the mission, the architecture and
the most serious technical risk. Conclude by defending one design compromise.
\end{speaking}
\end{multicols}

\subsection*{Sources}
\ECsource{NASA Small Spacecraft Systems Virtual Institute}{https://www.nasa.gov/smallsat-institute/}
\ECsource{European Space Agency -- Space Engineering and Technology}{https://www.esa.int/Enabling_Support/Space_Engineering_Technology}
